\chapter{Arquitetura da Plataforma}
\label{cap:arquitetura}

Este capítulo detalha a concepção arquitetural da plataforma educacional, demonstrando como a transição do protótipo inicial executado puramente no cliente para um modelo híbrido e reativo integrado cliente-servidor viabiliza avaliações seguras, supervisão pedagógica em tempo real, suporte síncrono em sala de aula e escalabilidade no ambiente acadêmico.

% ----------------------------------------------------------
\section{Visão Geral da Arquitetura Híbrida e Reativa}
\label{sec:visao-arquitetural}
% ----------------------------------------------------------

O projeto arquitetural da plataforma foi concebido para equacionar dois requisitos aparentemente antagônicos: por um lado, a necessidade de proporcionar uma experiência de terminal com \textbf{latência zero e alta responsividade} para o estudante; por outro lado, a exigência indispensável de \textbf{segurança avaliativa, integridade contra fraudes e supervisão docente contínua}.

No modelo convencional de aplicações web educacionais baseadas em requisições HTTP REST síncronas, a validação de tarefas práticas impõe atritos significativos ao usuário: o estudante precisa interromper seu fluxo de digitação no terminal, mudar o foco para a interface gráfica e acionar manualmente botões como ``Enviar Resposta'' ou ``Validar Tarefa''. Esse paradigma fragmenta a concentração do aprendiz e introduz latências de rede indesejadas a cada tentativa.

Para superar essa limitação, a plataforma adota uma **arquitetura reativa bidirecional orientada a eventos**, sustentada por canais persistentes seguros via \textit{WebSockets} (\textit{WSS}), operando em conformidade com o seguinte arranjo de responsabilidades:

\begin{itemize}
    \item \textbf{Camada de Execução Interativa (Cliente / Navegador):} O núcleo do simulador POSIX, a pilha de sessões Bash e a árvore do Sistema de Arquivos Virtual (\textit{VFS}) operam em memória local no navegador do aluno, desenvolvidos em TypeScript puro. Essa decisão assegura que a digitação de caracteres, a navegação pelo histórico, o autocompletar via tecla \texttt{Tab} e a renderização de cores ANSI ocorram instantaneamente, sem gerar tráfego de rede e sem sobrecarregar os servidores institucionais com a instanciação de contêineres individuais por aluno;
    \item \textbf{Camada de Validação e Orquestração (Servidor / Back-end):} O servidor atua como autoridade central confiável (\textit{Trusted Authority}). As regras de correção, os gabaritos das questões, os pesos das avaliações e as rotinas de asserção residem com exclusividade no servidor, tornando-os completamente inacessíveis à inspeção discente via console de desenvolvedor (\textit{DevTools});
    \item \textbf{Canal de Supervisão Docente (Cockpit do Professor):} O fluxo contínuo de eventos transmitidos via \textit{WebSockets} alimenta um painel administrativo em tempo real para o docente. A professora tem a capacidade de monitorar o progresso simultâneo de todos os terminais da turma durante uma prova prática presencial, identificar imediatamente discentes estagnados em mensagens de erro recorrentes e intervir pedagogicamente por meio de mensagens direcionadas ou avisos sincronizados.
\end{itemize}

% ----------------------------------------------------------
\section{Modelo de Dados Relacional}
\label{sec:modelo-dados}
% ----------------------------------------------------------

A gestão pedagógica e a persistência de longo prazo assentam-se sobre um banco de dados relacional estruturado (PostgreSQL), modelado para contemplar as entidades fundamentais do ecossistema acadêmico:
\begin{itemize}
    \item \textbf{Usuários e Perfis:} Diferenciação rigorosa entre os níveis de acesso Público, Aluno (vinculado por Registro Acadêmico - RA) e Professor;
    \item \textbf{Turmas e Inscrições:} Agrupamento semestral de turmas da disciplina de Sistemas Operacionais, possibilitando a matrícula e o acompanhamento de turmas regulares e especiais;
    \item \textbf{Avaliações e Janelas de Aplicação:} Modelagem de provas práticas associadas a janelas temporais estritas e chaves dinâmicas de liberação (\textit{tokens}), garantindo suporte nativo a repescagens e segundas chamadas isoladas;
    \item \textbf{Sessões e Submissões:} Registro auditável de cada tentativa individual de prova, armazenando o carimbo de data/hora validado pelo servidor, a nota obtida e a trilha completa de telemetria dos comandos executados.
\end{itemize}

% ----------------------------------------------------------
\section{Protocolo de Validação Desacoplada de VFS e Telemetria em Tempo Real}
\label{sec:validacao-desacoplada}
% ----------------------------------------------------------

Para mitigar integralmente os riscos de fraude via ferramentas de desenvolvedor (\textit{DevTools}) e, simultaneamente, eliminar a necessidade de botões manuais de envio que quebram o fluxo de raciocínio do discente, a validação de cada desafio opera através de um fluxo assíncrono de eventos via \textit{WebSockets}.

\subsection{Fluxo Reativo de Feedback Instantâneo}
A interação entre o terminal local e o motor de avaliação em retaguarda segue um ciclo não-bloqueante de três etapas:
\begin{enumerate}
    \item \textbf{Execução e Emissão em Background:} A cada quebra de linha (\texttt{Enter}) no terminal, o simulador executa a instrução instantaneamente no VFS local. Paralelamente à exibição imediata do resultado visual na tela do aluno, um emissor em segundo plano despacha pelo canal \textit{WebSocket} um \textit{payload} leve contendo o comando digitado, o código de retorno (\textit{exit code}), a identificação da questão e a árvore de nós ou metadados alterados (\textit{snapshot} diferencial);
    \item \textbf{Asserção em Ambiente Seguro:} O motor de avaliação no \textit{back-end} reconstitui o estado do VFS em memória protegida e aplica a bateria de testes de aceitação correspondente à tarefa ativa. Como as regras de teste residem unicamente no servidor, qualquer tentativa de manipulação de variáveis locais no navegador é inócua;
    \item \textbf{Retorno de Evento e Conclusão Automática:} Se as asserções forem plenamente satisfeitas, o servidor emite de volta pelo socket um evento assíncrono \texttt{task:completed}. A interface do aluno intercepta esse evento e exibe instantaneamente o \textit{feedback} visual de sucesso (animação de conclusão e avanço no roteiro), mantendo o aprendiz em estado contínuo de imersão e foco cognitivo.
\end{enumerate}

\subsection{Natureza Multifacetada do Linux e Auditoria Comportamental}
Essa abordagem fundamenta-se na própria natureza do ecossistema Linux, no qual inexiste um comando único para resolver determinada tarefa operacional. Por exemplo, a atribuição de permissões de leitura e escrita a um arquivo pode ser executada por notação octal (\texttt{chmod 644 arquivo.txt}) ou simbólica (\texttt{chmod u=rw,go=r arquivo.txt}); a busca por padrões de texto pode ser conduzida combinando \texttt{cat} e \texttt{grep} via \textit{pipes} ou diretamente através do argumento de entrada do utilitário \texttt{grep}.

Dessa forma, o motor de retaguarda atua em dois níveis complementares:
\begin{itemize}
    \item \textbf{Validação por Estado (Determinística):} O servidor avalia exclusivamente o resultado final refletido no VFS (verificando se o arquivo existe no caminho esperado, se possui o UID correto e se as permissões atendem à máscara exigida), assegurando aprovação pedagógica independentemente da rota sintática adotada pelo discente;
    \item \textbf{Análise Comportamental e Auditoria em Tempo Real:} A persistência cronológica de todas as instruções permite à docente mapear a proficiência dos comandos utilizados, diagnosticar lacunas conceituais recorrentes e garantir a integridade formal da avaliação contra tentativas de trapaça.
\end{itemize}

% ----------------------------------------------------------
\section{Governança Avaliativa, Telemetria de Segurança e Políticas Anti-Fraude}
\label{sec:governanca-seguranca}
% ----------------------------------------------------------

A aplicação de avaliações práticas em ambiente computacional exige um arcabouço rigoroso de governança que equilibre a supervisão pedagógica em tempo real, a prevenção ativa contra tentativas de trapaça e a tolerância a comportamentos acidentais. Em consonância com a arquitetura reativa orientada a eventos, a plataforma estrutura seus mecanismos de controle em cinco eixos operacionais: o suporte síncrono em sala de aula, a detecção heurística de ferramentas de inspeção, o bloqueio de concorrência de sessões, a telemetria forense de rede e a gestão de janelas temporais com chaves dinâmicas.

\subsection{Dinâmica Pedagógica da Fila de Suporte: A Mãozinha Virtual}
\label{sec:fila-suporte}

Em aulas laboratoriais práticas com turmas numerosas, um dos gargalos pedagógicos mais críticos reside na gestão do atendimento a dúvidas: múltiplos alunos erguem os braços simultaneamente, o docente perde a ordem cronológica dos chamados e dúvidas idênticas acabam sendo respondidas individualmente de forma exaustiva e repetitiva. Para solucionar esse problema, a plataforma incorpora o protocolo síncrono da \textbf{Fila de Suporte (Mãozinha Virtual)}, estruturado em três etapas operacionais integradas via \textit{WebSockets}:

\begin{enumerate}
    \item \textbf{Abertura Qualificada do Pedido de Ajuda (Levantar a Mão):} Ao acionar o comando de suporte na interface do terminal (``Pedir Ajuda''), o sistema exige obrigatoriamente que o estudante preencha um resumo textual sucinto descrevendo a natureza da sua dúvida ou o obstáculo enfrentado. Essa exigência pedagógica exercita a capacidade de síntese e metacognição do discente, forçando-o a estruturar o problema conceitualmente e eliminando chamados vazios ou acidentais. Uma vez despachado o evento \texttt{help:request}, o servidor insere um novo cartão na fila do painel docente (\textit{Cockpit}), ordenado pela disciplina cronológica FIFO (\textit{First-In, First-Out}), com data/hora, identificação do discente, terminal e descrição;
    \item \textbf{Cancelamento Autônomo (Abaixar a Mão):} Durante a espera na fila, caso o estudante consiga resolver o problema sozinho ou acompanhar uma explicação coletiva em sala, ele pode acionar a qualquer momento a opção ``Abaixar a Mão'', despachando \texttt{help:cancel}. O servidor remove instantaneamente o cartão da fila docente via \textit{WebSocket}, valorizando a autorregulação e evitando atendimentos desnecessários;
    \item \textbf{Atendimento Privado e Broadcast Inteligente:} A professora dispõe de dois modos de atendimento: a resposta individual discreta (diálogo 1:1 sobreposto ao terminal do aluno) e a resposta em lote por transmissão pública (\textit{multiselect}). Na resposta em lote, a docente seleciona múltiplos chamados correlatos e redige uma única orientação, a qual é transmitida para a turma mencionando os estudantes contemplados (exemplo: ``\textit{Respondendo às dúvidas de Aluno A, Aluno B e Aluno C: [Orientações da Docente]}''), resolvendo simultaneamente todos os chamados agrupados.
\end{enumerate}

\subsection{Políticas de Integridade e Detecção de DevTools com Avisos Escalonados}
\label{sec:integridade-devtools}

A integridade de uma avaliação prática computadorizada demanda proteção ativa contra a inspeção e a manipulação do código de execução no cliente. No entanto, em ambientes educacionais, o disparo de atalhos de teclado (como \texttt{F12}, \texttt{Ctrl+Shift+I} ou a alternância rápida de janelas com \texttt{Alt+Tab}) pode ocorrer de forma involuntária ou acidental pelo estudante. Para mitigar esse atrito sem prescindir do rigor avaliativo, o sistema implementa uma política de avisos escalonados e detecção em tempo real:

\begin{enumerate}
    \item \textbf{Sensoriamento Contínuo no Cliente:} O cliente mantém sentinelas em segundo plano que monitoram variações anômalas no redimensionamento da janela de exibição (\textit{window inner/outer dimensions mismatch}), limiares de tempo de execução causados por pontos de interrupção (\textit{debugger timing traps}) e eventos de perda de visibilidade da página por meio da API W3C \texttt{visibilitychange};
    \item \textbf{Primeira Ocorrência (Advertência Pedagógica):} Na primeira detecção de abertura de ferramentas de depuração ou perda de foco prolongada da aba de prova, o sistema não penaliza imediatamente a nota. Em vez disso, bloqueia temporariamente a digitação e apresenta um modal instrutivo em tela cheia, advertindo formalmente o discente sobre as regras de integridade e registrando a advertência local;
    \item \textbf{Reincidência e Notificação ao Painel Docente:} Caso o comportamento persista ou o estudante reabra as ferramentas de desenvolvedor, o cliente despacha imediatamente o evento \texttt{security:violation} pelo canal \textit{WebSocket}. O servidor registra a infração na trilha de telemetria forense e atualiza o \textit{card} do aluno no painel do professor com um indicador visual de alerta crítico (\textit{badge} de atenção). O sistema preserva a soberania da autoridade pedagógica, conferindo à docente a prerrogativa de verificar os registros circunstanciais, advertir presencialmente o aluno, desconsiderar falsos positivos ou anular a tentativa avaliativa.
\end{enumerate}

\subsection{Bloqueio de Concorrência de Sessão (Mutex Anti-Proxy)}
\label{sec:mutex-sessao}

Uma vulnerabilidade recorrente em exames computadorizados é o compartilhamento indevido de credenciais de acesso, no qual um discente presente no laboratório cede seu login a um colaborador externo (\textit{proxy} humano) para a resolução remota paralela da prova.

Para erradicar essa prática, a plataforma implementa uma trava de exclusão mútua de sessão (\textit{Session Mutex}) gerenciada no servidor sobre o protocolo de \textit{WebSockets}:
\begin{itemize}
    \item \textbf{Unicidade Estrita de Conexão:} Cada estudante autenticado por Registro Acadêmico (RA) só pode manter exatamente uma conexão \textit{socket} ativa por avaliação;
    \item \textbf{Derrubada Automática e Registro de Colisão:} Se uma nova tentativa de conexão for estabelecida com as mesmas credenciais (seja a partir de outro navegador, aba em modo anônimo ou dispositivo remoto distinto), o servidor derruba instantaneamente o canal de comunicação anterior com o código de encerramento \texttt{4409 Conflict} e emite a notificação: ``\textit{Conexão Encerrada: Nova sessão de avaliação detectada em outro dispositivo}'';
    \item \textbf{Alerta Imediato de Fraude:} O evento de colisão gera uma notificação prioritária no painel docente em tempo real, informando o endereço IP de ambas as pontas e o carimbo de data/hora da tentativa, permitindo a flagrância do compartilhamento indevido.
\end{itemize}

\subsection{Telemetria Forense de Rede, Geolocalização e Limitações de Sandbox}
\label{sec:telemetria-forense}

Como camada complementar de auditoria pós-exame, cada submissão e evento crítico de sessão é enriquecido com metadados de contexto de rede e ambiente operacional:
\begin{itemize}
    \item \textbf{Captura e Triagem de Endereços IP:} O servidor extrai o endereço IP público de origem e analisa o cabeçalho de proxy reverso corporativo. Em exames presenciais, o sistema permite à docente configurar um filtro de sub-rede institucional, confirmando automaticamente se todas as conexões provêm do bloco de rede e do \textit{gateway} exclusivo do laboratório físico da UTFPR;
    \item \textbf{Fingerprinting de Navegador e Coordenadas:} O cliente realiza uma composição de atributos de ambiente (\textit{User-Agent}, resolução de vídeo, idioma do sistema e \textit{timezone}), associando um \textit{hash} de integridade à sessão. Em avaliações autorizadas para modalidade remota, quando houver consentimento explícito do usuário, o sistema pode coletar as coordenadas aproximadas via API W3C de Geolocalização, subsidiando investigações de comissões acadêmicas em casos de suspeita fundamentada;
    \item \textbf{Delimitação de Sandbox e Privacidade W3C:} Cabe ressaltar formalmente que, devido às especificações universais de segurança e privacidade dos navegadores modernos mantidas pelo consórcio W3C, aplicações web executadas em modo \textit{sandbox} são estritamente impedidas de acessar identificadores de hardware de camada de enlace (como o endereço físico \textit{MAC Address} da interface de rede). Portanto, a eficácia do modelo de integridade apoia-se deliberadamente na combinação entre barreiras da camada de aplicação (tokens dinâmicos e travas de mutex) e verificações na camada de rede (triagem de IP institucional e canais criptografados TLS).
\end{itemize}

\subsection{Tokens Dinâmicos de Liberação e Janelas Temporais de Aplicação}
\label{sec:seguranca-prova}

O controle temporal e a circunscrição espacial dos exames completam o protocolo de governança da plataforma:
\begin{itemize}
    \item \textbf{Tokens Dinâmicos de Liberação em Sala:} No início da aula presencial de avaliação, a docente projeta ou divulga um código de sessão temporário com validade restrita ao horário letivo, impedindo que discentes iniciem a prova antecipadamente ou a realizem fora do laboratório físico;
    \item \textbf{Temporizador Sincronizado no Servidor:} A contagem regressiva e o encerramento do prazo de submissão baseiam-se estritamente no relógio oficial do servidor, inviabilizando fraudes por manipulação do relógio do sistema operacional do computador do aluno;
    \item \textbf{Janelas Isoladas para Repescagens e Segundas Chamadas:} Para discentes amparados por justificativas acadêmicas, o sistema permite o agendamento de janelas temporais dedicadas vinculadas exclusivamente aos seus respectivos RAs, liberando a avaliação em regime de reposição sem expor o exame à turma regular.
\end{itemize}
